\documentclass[11pt]{article}

\usepackage[margin=1in]{geometry}
\usepackage{amsmath,amssymb}
\usepackage{booktabs}
\usepackage{array}
\usepackage{graphicx}
\usepackage[round,authoryear]{natbib}
\usepackage[hidelinks]{hyperref}

\title{From Fixed to Evidence-Adaptive Persistence in Statistical Jump Models:\\
A Walk-Forward Three-Market Proxy Study}
\author{Tan Le}
\date{21 July 2026}

\newcommand{\BH}{Buy\&Hold}
\newcommand{\JM}{JM}
\newcommand{\HMM}{HMM}
\newcommand{\ind}{\mathbf{1}}
\newcommand{\F}{\mathcal{F}}
\newcommand{\sech}{\operatorname{sech}}
\setlength{\emergencystretch}{3em}

\begin{document}
\maketitle

\begin{abstract}
Financial markets often pass through persistent favorable and unfavorable
periods. This creates a simple investment problem: leaving equity during a bad
period can reduce downside risk, but an unstable signal can replace market risk
with late reactions and costly whipsaws. Shu, Yu, and Mulvey study this problem
with a transparent strategy: hold an equity index in a favorable regime and
local cash in an unfavorable regime. Their statistical Jump Model (JM)
encourages persistent states with a fixed transition penalty and, on their
exact 1990--2023 data, outperforms both buy-and-hold and a Gaussian hidden
Markov model (HMM).

We reconstruct the same past-only economic test on a later public proxy sample
through 2023 and ask whether the fixed switching barrier should respond to
evidence. We formulate and audit three time-varying decoders: arrival-day
evidence, lagged evidence, and a pair-balanced lagged penalty. They retain exact
dynamic programming, nest fixed JM at \(\beta=0\), and use a robust scale
estimated only from the training prefix. In the binary case, the
value-difference recursion shows why using the current loss gap twice can
amplify a shock; lagging the gap removes that direct double-use, while pair
balancing fixes the average two-direction transition scale.

Economically, the extension is not yet a general success. Fixed JM does not
beat the strongest same-sample benchmark in the United States, Germany, or
Japan. Even an ex-post upper envelope that chooses the best tested variant
separately in each market beats both benchmarks only in Germany. This envelope
is not one deployable model: it remains 0.037 Sharpe below HMM in the United
States and 0.131 below buy-and-hold in Japan. The contribution at this stage is
a candidate mathematical extension, a look-ahead-free audited experiment, and a clear
map of heterogeneous failure---not evidence of robust or generalizable
profitability. The evidence points toward a different next hypothesis: market
states may need to be formed partly by their matured future economic
consequences rather than by feature compactness and persistence alone.
\end{abstract}

\noindent\textbf{Keywords:} market regimes; statistical jump models; hidden
Markov models; downside risk; market timing; transition penalties.

\section{Introduction}

Traded markets do not behave as if one distribution governs every day.
Extended periods of rising or falling prices are commonly described as bull
and bear markets \citep{pagan2003}; volatility changes materially over time
\citep{schwert1989}; and investors distinguish risk-on from risk-off episodes
\citep{smales2016}. Regime-switching behavior has been studied in equities
\citep{hardy2001}, fixed income, currencies, and the business cycle.
Latent-state and turning-point methods offer different ways to date these
periods \citep{harding2003}. One reason the regime language persists is
interpretability: an inferred state can often
be related after the fact to a recession, crisis, recovery, or change in market
risk \citep{hamilton1989,ang2012}.

The investment problem is narrower than explaining every fluctuation. Losses
during a persistent unfavorable period can dominate long-horizon downside
risk. Safety-first and semivariance ideas already recognize that investors
need not value upside and downside symmetrically
\citep{roy1952,markowitz1959}. In regime-dependent allocation models, a
credible unfavorable state creates an economic reason to reduce risky
exposure and move toward a safe asset \citep{ang2002,ang2004}. The challenge is
timing. A signal that reacts to every noisy day may trade repeatedly and pay
costs; a signal that insists on too much persistence may recognize a crash only
after much of the loss has occurred. We call this the latency--whipsaw
trade-off.

A useful experiment should make the value of the signal visible. The 0/1
strategy does this by holding 100\% equity in a favorable state and 100\%
local cash in an unfavorable state. It deliberately removes portfolio
optimization and cross-asset diversification, so the comparison mostly asks
whether the regime signal itself is economically useful. HMM-guided versions
of this strategy provide an established benchmark \citep{bulla2011}. The
experiment is severe but easy to audit: a state becomes a signal, the signal
becomes a delayed position, and every position change pays a transaction cost.

An HMM and a JM express persistence differently. A Gaussian HMM assumes an
unobserved Markov chain and estimates state-transition probabilities and
conditional return distributions by likelihood. A fixed JM is closer to
clustering: it assigns observations to centers, but charges a constant cost
whenever two adjacent observations receive different states
\citep{bemporad2018}. This explicit jump penalty can suppress the short-lived
states that arise when a model follows daily noise. Earlier work connected
jump penalties to persistent market-state identification
\citep{nystrup2020persistent}; Shu, Yu, and Mulvey then asked whether that
statistical persistence produces a better delayed and costed investment
strategy \citep{shu2024}.

Shu et al.\ make two central moves. First, they evaluate persistence by its
economic consequence in a simple market/cash strategy rather than by a
clustering score alone. Second, they select the JM penalty each month with
time-series validation Sharpe after the same trading delay and cost used in
testing. On their licensed 1990--2023 sample, JM improves downside-risk and
risk-adjusted metrics relative to both HMM and \BH{} in the United States,
Germany, and Japan. Their paper is consequently not simply fitting JM and
comparing clusters. It is a walk-forward economic test of whether a more persistent
regime-identification signal is useful.

The fixed JM nevertheless treats every possible switch alike. A move supported
by an overwhelming loss difference between states pays the same \(\lambda\)
as a move supported by almost no difference. This observation leads to the
simple idea studied here: keep the fixed penalty as the default, but lower the
barrier when past-only state-fit evidence favors the destination. The formula
is simple; its consequences are not. The evidence can be counted twice,
directional costs can alter path boundaries, and monthly validation can select
different lambdas even when candidate paths improve mechanically.

We therefore ask two linked questions:
\begin{quote}
\textbf{Mathematical question.} Can a loss-gap-dependent switching cost be
inserted into JM without losing exact dynamic programming, training-only scaling, or
the fixed-JM special case?

\textbf{Economic question.} Can one prespecified JM variant beat both
same-sample \BH{} and Gaussian HMM after a one-day trading delay and 10-basis-
point one-way costs?
\end{quote}

Throughout this manuscript, ``causal'' refers only to information timing: each
decision uses quantities available by its decision date. It does not denote
treatment-effect identification.

The study makes three bounded contributions. First, it formulates a family of
past-only, loss-gap-dependent transition costs that retains
\(O(TK^2)\) dynamic programming and exactly nests fixed JM at \(\beta=0\).
Second, for two states it derives the value-difference recursion, identifies a
same-day evidence-amplification channel, and constructs lagged and
pair-balanced alternatives with an exact transition-cost identity. Third, it
documents heterogeneous economic effects on three public proxy development
samples: Germany improves substantially, but no tested variant beats both
benchmarks in every market. These statements do not establish global novelty,
statistical significance, or performance on untouched data.

\section{The Shu et al.\ parent study}

\subsection{Economic design}

Shu et al.\ use daily total-return indices for the S\&P 500, DAX, and Nikkei
225 from Bloomberg, together with each country's three-month Treasury-bill
yield from Global Financial Data. Their raw history begins in 1970. A
3,000-day training window and an eight-year validation period make the reported
testing period 1990--2023. The markets are tested independently in local
currency.

Their two-state HMM observes daily log total returns. It is refit daily on
3,000 observations from ten starts; the accepted fit has the highest
likelihood. The lower-volatility state is favorable and the higher-volatility
state is unfavorable. The terminal Viterbi state is an online state estimate.
Because that endpoint can be noisy, a causal trailing median filter is applied,
and its window length is selected monthly.

Their two-state JM observes three standardized excess-return features:
downside deviation with half-life 10 and Sortino ratios with half-lives 20 and
60. Centers are refit every six months on 3,000 observations from ten starts;
the fit with the lowest objective is retained. The state with higher cumulative
training excess return is labeled favorable. For each online date, the fitted
centers are frozen, the dynamic program is run on the trailing feature window,
and the terminal state is emitted. At the start of each month, the preceding
eight calendar years select the JM penalty and HMM smoothing value that
maximize delayed, costed validation Sharpe.

The state recognized after day \(t\) cannot earn day \(t+1\). Allocation is
changed at the end of \(t+1\), so the first affected return is \(t+2\).
Every one-way position change costs 10 basis points. This separates model
latency---the observations needed to recognize a new state---from operational
trading delay after the signal already exists.

\subsection{What Shu et al.\ report}

Table~\ref{tab:shu-main} reproduces the numerical contents of Shu et al.'s
primary performance table. Their ordering is consistent in all three markets:
JM has the highest Sharpe and Calmar ratios, lower maximum drawdown than both
comparators, and much lower turnover than HMM. Their claim is primarily about
downside-risk reduction and risk-adjusted performance, not maximum raw return.

\begin{table}[!htbp]
\centering
\caption{Results reported by Shu et al.\ for 1990--2023. Turnover is their
paper convention. These are parent-study values, not results produced by this
repository.}
\label{tab:shu-main}
\resizebox{\textwidth}{!}{
\begin{tabular}{llrrrrrrrr}
\toprule
Market & Strategy & Return & Vol. & Sharpe & MDD & Calmar & ES 5\% & Turnover & Equity exposure \\
\midrule
US & \BH{} & 10.2\% & 18.2\% & 0.48 & -55.2\% & 0.16 & -2.7\% & 0\% & 100\% \\
US & HMM    &  8.5\% & 11.3\% & 0.54 & -28.9\% & 0.21 & -1.8\% & 141\% & 72\% \\
US & JM     & 11.2\% & 13.1\% & 0.68 & -26.6\% & 0.33 & -2.0\% & 44\% & 80\% \\
\addlinespace
DE & \BH{} & 6.8\% & 22.1\% & 0.30 & -72.7\% & 0.09 & -3.3\% & 0\% & 100\% \\
DE & HMM    & 6.4\% & 14.0\% & 0.35 & -40.5\% & 0.12 & -2.2\% & 246\% & 73\% \\
DE & JM     & 8.6\% & 16.4\% & 0.44 & -39.4\% & 0.18 & -2.5\% & 170\% & 84\% \\
\addlinespace
JP & \BH{} & 0.8\% & 23.4\% & 0.12 & -79.1\% & 0.04 & -3.4\% & 0\% & 100\% \\
JP & HMM    & 2.5\% & 16.0\% & 0.19 & -48.6\% & 0.06 & -2.5\% & 290\% & 68\% \\
JP & JM     & 4.7\% & 17.1\% & 0.31 & -45.3\% & 0.12 & -2.6\% & 72\% & 75\% \\
\bottomrule
\end{tabular}}
\end{table}

Shu et al.\ also align validation and testing delays at 1, 5, and 10 trading
days. Their JM Sharpe ratios at those delays are \(0.68/0.71/0.70\) in the US,
\(0.44/0.38/0.29\) in Germany, and \(0.31/0.27/0.24\) in Japan. The comparison
supports their interpretation that a persistent state signal can remain useful
under longer execution delays.

\subsection{What the paper and public code do not disclose}

The authors' public repository contains a generic JM library and a Nasdaq
example. We audited the public code at commit
\texttt{d0fa00ce10126791695a259a45c5ddd41fbced80}. It confirms the half-squared
Euclidean loss, a zero diagonal and constant off-diagonal penalty matrix,
coordinate descent, dynamic programming, and ten-start fitting. It does not
contain the paper's Bloomberg/GFD data, HMM comparison, \BH{} path, monthly
validation, trading accounting, three-market replication, or final candidate
grids.

The paper explains the selection procedure but does not publish the complete
JM-lambda grid, HMM-smoothing grid, tie-breaking rule, exact semiannual refit
dates, raw vendor identifiers, cash-yield conversion, or machine-readable
states and trades. Table 3 displays illustrative fixed values
\(\lambda=\{0,5,15,35,70,150\}\) and
\(k=\{0,2,4,8,20\}\); it does not identify them as the complete validation
grids. Exact replication of Table~\ref{tab:shu-main} is therefore
underidentified without additional source data and conventions.

\section{Public proxy data}

We use public series capped at 31 December 2023. Table~\ref{tab:data} states
the definitions explicitly. These proxies were chosen to approximate the
economic roles in the parent study, not to claim vendor-level equivalence.

\begin{table}[!htbp]
\centering
\caption{Public proxy definitions and primary evaluation samples.}
\label{tab:data}
\small
\resizebox{\textwidth}{!}{%
\begin{tabular}{p{0.07\textwidth}p{0.21\textwidth}p{0.23\textwidth}p{0.30\textwidth}r}
\toprule
Market & Equity proxy & Cash proxy & Important difference from Shu et al. & Observations \\
\midrule
US & Yahoo \texttt{\char94 SP500TR}, gross total-return candidate
& FRED \texttt{DTB3}, three-month Treasury-bill discount rate
& Equity history begins in 1988, so the 1970 warm-up and 1990 test start are unavailable
& 4,046 \\
DE & Yahoo \texttt{\char94 GDAXI}, DAX performance-index candidate
& FRED \texttt{IR3TIB01DEM156N}, monthly interbank-rate proxy
& Equity history begins in 1987; cash is not a Treasury-bill series
& 4,059 \\
JP & Yahoo \texttt{\char94 N225}, Nikkei 225 price index
& Bank of Japan \texttt{STRACLUC3M}, monthly uncollateralized call-rate proxy
& The equity series excludes dividends and cash is not a Treasury-bill series
& 3,586 \\
\bottomrule
\end{tabular}
}
\end{table}

Equity simple return is \(P_t/P_{t^-}-1\), where \(t^-\) is the preceding
valid equity date; log return is \(\log(P_t/P_{t^-})\). Annual percentage cash
rates are converted to daily returns by division by \(100\times252\) and are
aligned backward only after causal availability. Excess return is equity
simple return minus cash return. Missing source values are preserved rather
than filled from the future.

After feature warm-up, selection, delay, and common-date alignment, the primary
samples are:

\begin{center}
\begin{tabular}{lccc}
\toprule
Market & Start & End & Observations \\
\midrule
US & 4 Dec 2007 & 29 Dec 2023 & 4,046 \\
DE & 3 Jan 2008 & 29 Dec 2023 & 4,059 \\
JP & 7 May 2009 & 29 Dec 2023 & 3,586 \\
\bottomrule
\end{tabular}
\end{center}

The periods are economically different from 1990--2023. The Japanese
difference is especially material: proxy \BH{} Sharpe is about 0.545, versus
0.12 in Shu et al., and the proxy omits dividends. Parent-paper performance is
therefore context. The scientific comparison here is the ordering of models on
the same proxy dates, not numerical equality with Table~\ref{tab:shu-main}.

\section{From an observed return to a tradable position}

Let \(r^e_t\) and \(r^c_t\) denote equity and local-cash returns. A model emits
\(a_t\in\{0,1\}\), with one denoting equity. The signal formed after observing
day \(t\) is executed at the end of the next trading day and first earns the
return at \(t+2\):
\begin{equation}
r^p_{t+2}=a_t r^e_{t+2}+(1-a_t)r^c_{t+2}.
\label{eq:gross}
\end{equation}
If \(w_u\) is the position that earns return \(u\), net return is
\begin{equation}
r^{p,\mathrm{net}}_u
=w_u r^e_u+(1-w_u)r^c_u
-0.001\,|w_u-w_{u-1}|.
\label{eq:net}
\end{equation}
The initial position does not receive a fictional trade cost.

\begin{figure}[!htbp]
\centering
\fbox{\parbox{0.92\textwidth}{
\centering
observe returns/features at \(t\)
\(\longrightarrow\) infer state \(s_t\)
\(\longrightarrow\) set signal \(a_t=1-s_t\)
\(\longrightarrow\) execute at end of \(t+1\)
\(\longrightarrow\) earn net return at \(t+2\)
}}
\caption{The common causal timeline. Every candidate, HMM, and fixed/adaptive
JM path uses this timing.}
\label{fig:timeline}
\end{figure}

This is regime identification rather than direct next-day return forecasting.
The model describes the state at \(t\); the strategy earns money only if the
state's risk--return relation persists after execution. Shu et al.\ explicitly
make this distinction, and related allocation work also separates
identification from forecasting \citep{nystrup2015}.

Annualized Sharpe is
\begin{equation}
\operatorname{Sharpe}
=\sqrt{252}\,
\frac{\operatorname{mean}(r^{p,\mathrm{net}}_t-r^c_t)}
     {\operatorname{sd}(r^{p,\mathrm{net}}_t)},
\label{eq:sharpe}
\end{equation}
using sample standard deviation. Maximum drawdown is the minimum peak-to-trough
wealth decline. Cash fraction is \(\operatorname{mean}(1-w_t)\).
Switch count is the number of nonzero adjacent position changes.

Turnover follows the parent-paper convention:
\begin{equation}
\operatorname{Turnover}
=\frac{1}{2}\,252\,
\operatorname{mean}|w_t-w_{t-1}|.
\label{eq:turnover}
\end{equation}
For a binary strategy, an equity-to-cash-to-equity round trip counts as one
unit. The old repository display omitted the one-half and therefore reported
annualized traded notional instead of paper turnover. Equation
\eqref{eq:net} always used the full one-way change, so correcting the turnover
label changed no cost, return, Sharpe, or drawdown.

\section{Benchmark regime models}

\subsection{Gaussian HMM}

An HMM assumes an unobserved two-state Markov chain \(s_t\) with transition
matrix \(P\), and a Gaussian distribution for the observed log return
conditional on each state. In simple language, it fits two return
distributions and probabilities of remaining in or moving between them. The
Baum--Welch expectation--maximization procedure estimates parameters by
likelihood; Viterbi decoding finds a most likely state path. Gaussian HMMs can
reproduce several stylized properties of daily equity returns
\citep{ryden1998}, but likelihood is multimodal and the assumed distribution
can be misspecified.

The proxy protocol fits a two-state diagonal-covariance Gaussian HMM daily on
the trailing 3,000 finite log returns. Ten deterministic starts are attempted
and the highest-likelihood accepted fit is kept. The state with lower fitted
variance is favorable. The terminal Viterbi state is recorded as the day's
raw online state. For each smoothing candidate \(k\), a causal trailing
majority/median filter acts only on states already emitted. \(k=0\) is
unsmoothed. The monthly validation procedure chooses \(k\).

\subsection{Fixed statistical Jump Model}

Let \(x_t\in\mathbb{R}^3\) be a standardized feature vector, let
\(\theta_0,\theta_1\) be centers, and let \(s_t\in\{0,1\}\). Fixed JM solves
\begin{equation}
\min_{\Theta,S}
\sum_{t=0}^{T}
\frac{1}{2}\|x_t-\theta_{s_t}\|_2^2
+\lambda\sum_{t=1}^{T}\ind\{s_t\ne s_{t-1}\}.
\label{eq:fixed}
\end{equation}
The first term is k-means-like compactness. The second charges \(\lambda\)
each time the path changes state. At \(\lambda=0\), time order does not affect
the assignment step; larger \(\lambda\) creates persistence, and an
excessively large value can collapse the path to one state.

The three v7 features follow the parent design:

\begin{center}
\begin{tabular}{lll}
\toprule
Feature & Half-life & Economic role \\
\midrule
Exponentially weighted downside deviation & 10 days & recent downside risk \\
Exponentially weighted Sortino ratio & 20 days & medium-horizon return/downside risk \\
Exponentially weighted Sortino ratio & 60 days & slower return/downside risk \\
\bottomrule
\end{tabular}
\end{center}

The learner alternates two simple steps. Holding the states fixed, it updates
each center from the observations assigned to that state. Holding the centers
fixed, it obtains the globally optimal path for those centers by dynamic
programming. Ten initializations are attempted and the lowest objective is
kept. Features are standardized from past data only. The proxy centers are
refit in January and July on a rolling 3,000-observation window. The state
with higher cumulative training excess return is favorable; the other is
cash.

For online date \(t\), centers remain frozen at the most recent refit. The
decoder sees the trailing 3,000 feature rows through \(t\), and its terminal
state becomes the signal state for \(t\). A large \(\lambda\) reduces switches
but can delay a turning point; a small \(\lambda\) reacts sooner but can
whipsaw. This is the same accuracy--latency problem emphasized in persistent
online classification \citep{nystrup2018,nystrup2020online}.

\subsection{Monthly hyperparameter selection}

At each prior-month final complete date, every eligible JM lambda and HMM
smoothing candidate is replayed over the preceding eight calendar years. Its
0/1 strategy uses the same \(t+2\) timing and 10-basis-point costs as the outer
sample. The candidate with highest validation Sharpe is activated for the
next month, subject to the same delay. Candidates with fewer than 252 valid
returns are ineligible. Ties within \(10^{-12}\) choose the smaller smoothing
value. Thus hyperparameter selection itself is causal and economic rather than
a retrospective choice from outer Sharpe.

\section{Evidence-adaptive transition penalties}

\subsection{General objective and dynamic program}

For fitted centers, define the state-fit loss
\begin{equation}
L_t(k)=\frac{1}{2}\|x_t-\theta_k\|_2^2.
\label{eq:loss}
\end{equation}
A time-varying decoder solves
\begin{equation}
\min_S
\sum_{t=0}^{T}L_t(s_t)
+\sum_{t=1}^{T}C_t(s_{t-1},s_t),
\qquad C_t(i,i)=0.
\label{eq:tvobjective}
\end{equation}
The exact recurrence is
\begin{equation}
V_t(j)=L_t(j)+\min_i\{V_{t-1}(i)+C_t(i,j)\}.
\label{eq:dp}
\end{equation}
Backtracking gives the minimizing path. The complexity remains
\(O(TK^2)\); for two states it is linear in \(T\).

At each refit,
\begin{equation}
q_{\mathrm{train}}
=\operatorname{median}_{u,k}
\left|L_u(k)-\operatorname{median}_{v,\ell}L_v(\ell)\right|
\label{eq:q}
\end{equation}
is the raw median absolute deviation of all finite state losses on the causal
3,000-row training prefix. It must be finite and strictly positive. This
training-only scale makes the loss gap dimensionless and prevents an asset's
feature scale from silently redefining \(\beta\). Adaptive studies reuse the
sealed v7 standardizer and fixed-JM centers for each lambda/refit; they change
online decoding only and do not jointly refit centers under
\eqref{eq:tvobjective}.

\subsection{Arrival-day evidence}

The first rule lowers the cost from previous state \(i\) to destination \(j\)
only when the current observation fits \(j\) better:
\begin{equation}
C^A_t(i,j)=
\lambda_0\exp\!\left[
-\beta\tanh\!\left(
\frac{[L_t(i)-L_t(j)]_+}{q_{\mathrm{train}}}
\right)\right],
\qquad i\ne j.
\label{eq:arrival}
\end{equation}
For \(\beta\ge0\),
\begin{equation}
\lambda_0e^{-\beta}\le C^A_t(i,j)\le\lambda_0.
\end{equation}
Thus \(\beta=\log2\) permits at most a 50\% discount and
\(\beta=\log4\) at most a 75\% discount. At \(\beta=0\), every off-diagonal
cost is exactly \(\lambda_0\), so the objective, states, choices, positions,
trades, costs, and returns equal fixed JM.

The intuitive weakness is same-day double-use. The current gap first enters
the observation loss and then lowers the current transition barrier. In the
binary case, let
\[
d_t=V_t(1)-V_t(0),\quad
g_t=L_t(1)-L_t(0),\quad
a_t=C_t(0,1),\quad b_t=C_t(1,0).
\]
Equation~\eqref{eq:dp} becomes
\begin{equation}
d_t=g_t+\operatorname{clip}(d_{t-1},-b_t,a_t).
\label{eq:binary}
\end{equation}
On the evidence-supported active boundary under arrival adaptation,
\begin{equation}
\left|\frac{\partial d_t}{\partial g_t}\right|
=
1+\frac{\beta C_t}{q_{\mathrm{train}}}
\sech^2\!\left(\frac{|g_t|}{q_{\mathrm{train}}}\right)\ge1.
\label{eq:amplification}
\end{equation}
The inequality is strict when \(\beta>0\), \(\lambda_0>0\), the evidence gap
is strict, and the relevant clip boundary is active away from kinks and ties.
At \(\beta=0\) or \(\lambda_0=0\), the derivative equals one.
This is not a general instability theorem, but it identifies a precise branch
where the same shock moves the loss difference and simultaneously weakens the
barrier.

\subsection{Lagged evidence}

The simplest repair changes only the evidence date:
\begin{equation}
C^L_t(i,j)=
\lambda_0\exp\!\left[
-\beta\tanh\!\left(
\frac{[L_{t-1}(i)-L_{t-1}(j)]_+}{q_{\mathrm{train}}}
\right)\right],
\qquad i\ne j.
\label{eq:lagged}
\end{equation}
The current observation still contributes \(L_t(s_t)\), but the previous
observation controls the barrier. This removes direct current-day double-use
while preserving the bounds, scale, and beta-zero identity. On January and
July refit dates, the previous row's loss is recomputed under the scaler and
centers fitted through \(t\). The rule is causal at \(t\), but it is not
strictly \(\F_{t-1}\)-measurable on those refit dates. Reusing stale centers
only for the evidence term would define a different \(\beta>0\) model, not the
frozen mechanism tested here. Beta-zero nesting itself is unaffected because
the loss gap then has no effect on the transition cost.

\subsection{Directed costs and pair balance}

An early project idea assigned constant costs \(a=C(0,1)\) and
\(b=C(1,0)\) and interpreted them as separate state durations. For a binary
path,
\begin{equation}
aN_{01}+bN_{10}
=\frac{a+b}{2}N_{\mathrm{switch}}
+\frac{a-b}{2}(s_T-s_0).
\label{eq:directed}
\end{equation}
A constant asymmetric matrix is therefore a symmetric switch cost plus an
endpoint boundary bias, not independent bull/bear duration control. That early
interpretation was withdrawn. When costs vary with time, the per-transition
decomposition remains valid, but the antisymmetric terms need not telescope.

The pair-balanced lagged rule is
\begin{equation}
C^B_t(i,j)=\lambda_0\left[
1-(1-e^{-\beta})
\tanh\!\left(
\frac{L_{t-1}(i)-L_{t-1}(j)}
     {q_{\mathrm{train}}}
\right)
\right],
\qquad i\ne j.
\label{eq:balanced}
\end{equation}
It satisfies
\begin{equation}
C^B_t(i,j)+C^B_t(j,i)=2\lambda_0
\label{eq:pairsum}
\end{equation}
at every date. Evidence tilts direction while the pair-average cost remains
fixed. Its bounds are
\(\lambda_0e^{-\beta}\le C^B_t(i,j)
\le\lambda_0(2-e^{-\beta})\).
It also nests fixed JM exactly at \(\beta=0\).

\subsection{A separation gate that did not survive testing}

The contrast between Germany and Japan suggested a simple hypothesis:
evidence discounts may work only when the two fitted states are well separated.
A training-only reliability statistic was frozen:
\begin{equation}
D=\|\mu_0-\mu_1\|,\qquad
\rho_k=\operatorname{median}_{u\in A_k}\|z_u-\mu_k\|,\qquad
R_{\mathrm{train}}=\frac{D}{D+\rho_0+\rho_1},
\label{eq:separation}
\end{equation}
where \(A_k\) contains training rows strictly nearer center \(k\). The proposed
gate multiplied \(\beta\) by \(R_{\mathrm{train}}\). It is bounded,
label-symmetric, and invariant to translation, orthogonal rotations, and
common positive scaling.

The completed performance-free diagnostic was inconclusive. All
leave-one-market-out optimizer fits failed the frozen gradient criterion, and
event reliability did not consistently distinguish whipsaws from persistent
moves. Germany even had slightly higher median reliability for whipsaw events
than persistent events. This gate was closed rather than promoted to a P\&L
sweep. It remains part of the mathematical record because it shows how a
plausible story can fail before trading performance is opened.

\section{Empirical design and verification}

\subsection{Economic target}

For market \(m\) and a prespecified JM variant \(v\), define
\begin{equation}
G_m(v)=
\operatorname{Sharpe}_{v,m}
-\max\{\operatorname{Sharpe}_{\mathrm{BH},m},
       \operatorname{Sharpe}_{\mathrm{HMM},m}\}.
\label{eq:gap}
\end{equation}
A market passes only if \(G_m(v)>0\). The project target requires the same
prespecified variant to pass all declared markets. The descriptive quantity
\(\max_v G_m(v)\), which chooses a different variant after seeing each
market, is an upper envelope and not a deployable model.

MDD, turnover, cash fraction, and switch count explain the risk and activity
used to obtain Sharpe; they do not replace the benchmark test. A switch count
is also not a true whipsaw label. Mechanism studies use an explicit 20-day
reversal definition, still without latent-regime ground truth.

\subsection{Frozen local choices}

The canonical proxy JM grid is
\[
\{0,5,15,35,70,150,300,600,1200\}.
\]
The HMM smoothing grid is
\[
\{0,2,4,6,8,10,20,40,80,160,320,640,1280,2560\}.
\]
These are repository choices, not recovered Shu parameters. Arrival and
lagged mechanism scenarios use
\(\beta\in\{0,\log2,\log4\}\). The lagged P\&L readout carries forward only
\(\log4\), selected by a performance-free mechanism rule. The later balanced
study compares \(\{0,\log4\}\), again without a new beta sweep. Beta is
scenario-fixed and not estimated to maximize profit.

A local upper-candidate concentration diagnostic was introduced only to warn
that a grid may be too short when the largest candidate is chosen in more than
5\% of months. It is not in Shu et al. and never changes a state, trade, cost,
or return. In fail-closed grid and window audits, it historically governed
whether performance metrics were opened. In accepted completed P\&L studies it
is descriptive only and does not hide reported metrics. An endpoint audit
later confirmed that lambda-domain truncation is real but did not rescue the
three-market ordering.

\subsection{Implementation contracts}

Before interpreting performance, the following contracts were checked:

\begin{table}[!htbp]
\centering
\caption{Verification contracts applied to the adaptive decoders and
accounting path.}
\label{tab:verification}
\small
\begin{tabular}{p{0.26\textwidth}p{0.64\textwidth}}
\toprule
Contract & Observable check \\
\midrule
Formula values & Stored penalties equal the stated equations, directions, bounds, and training-only scale \\
Objective & Reported path objective equals observation losses plus the exact time-indexed transition costs \\
Beta-zero nesting & Candidate states and full selected trade paths exactly match fixed JM at \(\beta=0\) \\
Brute-force equivalence & Dynamic programming equals exhaustive enumeration on small binary paths \\
Toy paths & Weak evidence retains the state; strong destination evidence can trigger the intended transition \\
Prefix invariance & Truncating future rows leaves all earlier states unchanged; future mutation cannot change the prefix \\
Causal fit boundaries & Each scaler, center, scale, candidate choice, signal, and position uses only information available by its decision date \\
Trading identity & Signal at \(t\) equals \(1-s_t\), position earns \(t+2\), and each one-way change costs exactly 0.001 \\
Concrete dates & Loss \(\rightarrow\) penalty \(\rightarrow\) state \(\rightarrow\) signal \(\rightarrow\) position \(\rightarrow\) trade is reconstructed \\
Independent replay & Accepted summaries, choices, positions, costs, and decisions are rebuilt from source artifacts \\
\bottomrule
\end{tabular}
\end{table}

The final balanced run independently reconstructed nine metric rows, 15,228
selection-surface rows, and 35,073 timeline rows with maximum numerical error
0.0. Focused tests passed 89 checks; the final repository suite passed 531.
These facts establish implementation health. They do not establish that an
inferred state is the market's unique true regime or that a strategy will
generalize.

\subsection{Development-sample status}

All US, DE, and JP proxy rows through 2023 have now been inspected repeatedly
during model development. The adaptive results are therefore exploratory
development evidence, even when a spec was frozen before a particular readout.
No claim of an untouched holdout, alpha, robustness, statistical significance,
stable profit, or live performance is authorized. No reported model or P\&L
experiment accessed post-2023 rows. A separate authorized source audit did
inspect public candidate series through July 2026, so those dates are not
untouched confirmation evidence.

\section{Results}

\subsection{The fixed parent ordering is not recovered}

Table~\ref{tab:paper-proxy} puts the central Sharpe comparison in one place.
The public proxy does not recover the parent ordering. Fixed JM is below the
strongest same-sample benchmark in all three markets.

\begin{table}[!htbp]
\centering
\caption{Paper Sharpe versus the current proxy fixed baseline. The columns are
not like-for-like replications because sources and evaluation periods differ.}
\label{tab:paper-proxy}
\begin{tabular}{lcc}
\toprule
Market & Shu et al.\ \BH{}/HMM/JM & Proxy \BH{}/HMM/fixed JM \\
\midrule
US & \(0.48/0.54/0.68\) & \(0.512964/0.653725/0.569865\) \\
DE & \(0.30/0.35/0.44\) & \(0.289638/0.007599/0.166440\) \\
JP & \(0.12/0.19/0.31\) & \(0.544589/0.398539/0.329270\) \\
\bottomrule
\end{tabular}
\end{table}

This is a proxy non-replication, not a refutation of Shu et al. Exact vendor
series, the 1970 warm-up, final grids, and several conventions are unavailable.
The Japanese proxy describes a particularly different investment: its price
index excludes dividends and its outer sample begins after the global
financial crisis.

\subsection{Full primary-delay proxy results}

Table~\ref{tab:all-results} reports all accepted primary-delay variants.
Arrival, lagged, and balanced mean \(\beta=\log4\). Turnover uses the corrected
paper convention. Values are rounded for display; the repository CSVs retain
full precision.

\begin{table}[p]
\centering
\caption{Primary-delay net results on the public proxy development sample
through 2023. Return, volatility, MDD, turnover, and cash are percentages.}
\label{tab:all-results}
\resizebox{\textwidth}{!}{
\begin{tabular}{llrrrrrrr}
\toprule
Market & Model & CAGR & Vol. & Sharpe & MDD & Turnover & Cash & Switches \\
\midrule
US & \BH{}       & 9.780 & 20.495 & 0.512964 & -53.963 & 0.000 & 0.000 & 0 \\
US & HMM          & 8.461 & 12.143 & \textbf{0.653725} & -19.363 & 143.253 & 24.221 & 46 \\
US & Fixed JM     & 8.512 & 14.595 & 0.569865 & -33.857 & 65.398 & 21.157 & 21 \\
US & Arrival JM   & 7.297 & 14.223 & 0.501745 & -33.790 & 59.170 & 23.900 & 19 \\
US & Lagged JM    & 8.508 & 14.026 & 0.586777 & -33.790 & 46.713 & 24.913 & 15 \\
US & Balanced JM  & 8.870 & 13.856 & 0.616326 & -33.790 & 40.484 & 26.174 & 13 \\
\addlinespace
DE & \BH{}       & 4.737 & 22.094 & 0.289638 & -53.876 & 0.000 & 0.000 & 0 \\
DE & HMM          & -1.031 & 19.240 & 0.007599 & -54.590 & 105.543 & 11.949 & 34 \\
DE & Fixed JM     & 2.071 & 17.118 & 0.166440 & -38.779 & 99.335 & 18.280 & 32 \\
DE & Arrival JM   & 4.123 & 18.082 & 0.277070 & -38.779 & 43.459 & 13.107 & 14 \\
DE & Lagged JM    & 5.259 & 17.987 & \textbf{0.337888} & -38.779 & 24.834 & 13.427 & 8 \\
DE & Balanced JM  & 5.098 & 17.291 & 0.335543 & -38.779 & 49.667 & 16.876 & 16 \\
\addlinespace
JP & \BH{}       & 9.687 & 20.731 & \textbf{0.544589} & -31.799 & 0.000 & 0.000 & 0 \\
JP & HMM          & 6.075 & 19.127 & 0.398539 & -34.656 & 130.006 & 5.633 & 37 \\
JP & Fixed JM     & 4.498 & 17.910 & 0.329270 & -32.162 & 45.678 & 27.886 & 13 \\
JP & Arrival JM   & 5.321 & 16.880 & 0.385254 & -41.348 & 94.869 & 27.524 & 27 \\
JP & Lagged JM    & 5.792 & 16.772 & 0.413215 & -31.799 & 115.951 & 27.970 & 33 \\
JP & Balanced JM  & 2.913 & 17.743 & 0.244542 & -41.303 & 137.033 & 22.644 & 39 \\
\bottomrule
\end{tabular}}
\end{table}

The fixed-JM result differs by market. In the US it lies between \BH{} and HMM.
In Germany it beats HMM but not \BH{}. In Japan it loses to both. Hence fixed
JM passes the economic target in \(0/3\) markets.

\subsection{Arrival adaptation}

The arrival study tested the complete frozen beta set. Sharpe at
\(\beta=0,\log2,\log4\) was:

\begin{center}
\begin{tabular}{lrrr}
\toprule
Market & \(\beta=0\) & \(\beta=\log2\) & \(\beta=\log4\) \\
\midrule
US & 0.569865 & 0.512427 & 0.501745 \\
DE & 0.166440 & 0.271817 & 0.277070 \\
JP & 0.329270 & 0.143204 & 0.385254 \\
\bottomrule
\end{tabular}
\end{center}

Beta zero reproduced fixed JM exactly. The mechanism changed paths and reduced
German switches from 32 to 14, so it behaved nontrivially. Economically,
however, arrival \(\log4\) beat both benchmarks in \(0/3\) markets. It lowered
US Sharpe and doubled Japanese activity from 13 to 27 switches while worsening
Japanese MDD by about 9.19 percentage points. Equation
\eqref{eq:amplification} gives a mathematical reason to suspect the current-day
double-use, but it does not by itself prove that double-use caused the market
result.

\subsection{Lagging the evidence}

Before P\&L was opened, a performance-free mechanism study compared
arrival-day and lagged candidate paths. At \(\beta=\log4\), 20-day whipsaw
events fell from \(6\) to \(2\) in the US, \(6\) to \(3\) in Germany, and
\(5\) to \(1\) in Japan: \(17\rightarrow6\) pooled. Across all three markets,
eleven persistent events still moved before fixed JM; Japan contributed five.
Only \(\log4\) passed the frozen mechanism
rule. This supported a mechanism experiment, not a return claim.

In the subsequent readout, lagged-minus-fixed Sharpe was
\(+0.016913/+0.171448/+0.083945\) in US/DE/JP. Switches changed
\(21\rightarrow15\), \(32\rightarrow8\), and \(13\rightarrow33\).
The model therefore reduced selected-strategy activity in the US and Germany
but increased it sharply in Japan. It beat both economic benchmarks only in
Germany. The apparent contradiction with the candidate-path whipsaw diagnostic
is resolved by monthly selection: the selector can choose different lambdas
when the entire candidate family changes.

A frozen \(2\times2\) diagnostic separated fixed/lagged candidate paths from
fixed/lagged monthly choices. For equal-market mean Sharpe, the total
lagged-minus-fixed change was \(+0.090769\). The direct path effect at fixed
choices was \(-0.137373\), the direct choice effect at fixed paths was
\(+0.023065\), and the interaction was \(+0.205076\). Shapley accounting
allocated \(-0.034835\) to the path family and \(+0.125603\) to the choice
schedule. These are exact decompositions of nonlinear metrics, not causal
effects. They show that the positive aggregate readout is not evidence that
lagged paths alone were uniformly better.

\subsection{Pair-balanced lagging}

The pair-balanced mechanism retained 87.5\% of eligible early confirmations
and produced zero unconfirmed persistent lock-ins, but pooled whipsaws did not
fall: own-event whipsaws were \(7\) versus \(6\) for one-sided lagging, and
matched-anchor whipsaws were \(5\) versus \(5\). The mechanism criterion was
therefore not supported.

The later, explicitly post-result P\&L readout found balanced-minus-lagged
Sharpe of \(+0.029549/-0.002345/-0.168672\) in US/DE/JP. US turnover and
switches fell; German turnover and switches rose; Japan deteriorated on Sharpe,
MDD, turnover, and switch count. Balanced-minus-fixed mean Sharpe was positive
with two positive markets, but the primary comparison with lagged had negative
mean Sharpe and only one positive market. Pair balancing was not supported as
the next model.

\subsection{Final benchmark gaps}

Table~\ref{tab:gaps} reports the most favorable observed JM separately in each
market. This is deliberately an ex-post upper envelope.

\begin{table}[!htbp]
\centering
\caption{Best observed per-market JM and gap to the strongest same-sample
benchmark. A single deployable variant cannot choose this row after seeing the
market.}
\label{tab:gaps}
\begin{tabular}{llrrrrrr}
\toprule
Market & Best JM & Sharpe & \(G_m\) & MDD & Turnover & Cash & Switches \\
\midrule
US & Balanced & 0.616326 & -0.037398 & -0.337905 & 0.404844 & 0.261740 & 13 \\
DE & Lagged   & 0.337888 & +0.048251 & -0.387794 & 0.248337 & 0.134270 & 8 \\
JP & Lagged   & 0.413215 & -0.131375 & -0.317989 & 1.159509 & 0.279699 & 33 \\
\bottomrule
\end{tabular}
\end{table}

Only Germany clears both \BH{} and HMM. The US remains below HMM; Japan remains
well below \BH{}. Fixed and arrival pass \(0/3\) markets; lagged and balanced
each pass only Germany. No tested variant meets the cross-market objective.

\subsection{Trading-delay sensitivity of the fixed proxy}

Table~\ref{tab:delay} reports the already accepted fixed-baseline delay study.
Each delay is applied consistently to validation and testing. Samples shorten
slightly as the offset grows. Adaptive variants were not tested at delays 5
and 10, so this table cannot support adaptive-delay robustness.

\begin{table}[!htbp]
\centering
\caption{Proxy Sharpe at trading delays 1/5/10 days.}
\label{tab:delay}
\begin{tabular}{llccc}
\toprule
Market & Model & Delay 1 & Delay 5 & Delay 10 \\
\midrule
US & \BH{} & 0.513 & 0.507 & 0.515 \\
US & HMM & 0.654 & 0.741 & 0.479 \\
US & Fixed JM & 0.570 & 0.548 & 0.604 \\
\addlinespace
DE & \BH{} & 0.290 & 0.294 & 0.304 \\
DE & HMM & 0.008 & 0.017 & 0.166 \\
DE & Fixed JM & 0.166 & 0.217 & 0.330 \\
\addlinespace
JP & \BH{} & 0.545 & 0.534 & 0.535 \\
JP & HMM & 0.399 & 0.397 & 0.404 \\
JP & Fixed JM & 0.329 & 0.256 & 0.315 \\
\bottomrule
\end{tabular}
\end{table}

The paths are not monotonically worse with longer delay because the validation
selector is itself re-optimized for each delay and the comparison dates change
slightly. Germany's fixed JM exceeds both comparators at delay 10 on that
shorter matched sample, whereas Japan never does. These heterogeneous values
do not reproduce Shu et al.'s broad delay-robustness pattern.

\subsection{Two concrete loss-to-trade traces}

Aggregate metrics can hide what the mechanism actually does.
Table~\ref{tab:traces} follows two accepted rows through the entire causal
chain. State 0 is favorable/equity and state 1 is unfavorable/cash.

\begin{table}[!htbp]
\centering
\caption{Concrete audited examples. The trade executes at the end of \(t+1\);
the displayed date is the \(t+2\) position/return row on which the one-way cost
is booked.}
\label{tab:traces}
\resizebox{\textwidth}{!}{
\begin{tabular}{llrrrrllllr}
\toprule
Market/rule & Evidence date & \(\lambda_0\) & \(q\) & Loss source & Loss dest. &
Penalty & State & Signal/position/trade & \(t+2\) position/return & Net return \\
\midrule
DE arrival & 25 Jan 2008 & 150 & 0.9901 & 7.4160 & 1.1158 &
37.5003 & \(0\to1\) & \(0/0/-1\) & 29 Jan 2008 & -0.000816 \\
JP lagged & 7 May 2009 & 5 & 0.9414 & 2.1566 & 0.8314 &
1.4620 & \(1\to0\) on 8 May & \(1/1/+1\) & 12 May 2009 & -0.017226 \\
\bottomrule
\end{tabular}}
\end{table}

In Germany, arrival evidence cut a potential \(0\to1\) transition cost from
150 to about 37.5; the state changed on 25 January. The trade executed at the
end of 28 January; the cash position first earned a return and booked the 0.001
one-way cost on 29 January. Fixed JM remained favorable on the signal date. In
Japan, the 7 May lagged loss gap cut the next day's \(1\to0\) cost from 5 to
about 1.462. The state changed to equity on 8 May and the trade executed at the
end of 11 May; the equity position first earned a return and booked the 0.001
cost on 12 May, when it lost roughly 1.72\% net. One date cannot
explain a full-sample Sharpe, but it makes the model and accounting observable.

\section{Discussion}

\subsection{What has been contributed mathematically}

The study has an explicit time-varying JM objective, a training-only robust
scale, exact dynamic programming, and an exact fixed-JM special case. For two
states, Equation~\eqref{eq:binary} compresses the decoder to a scalar
value-difference recursion. Equation~\eqref{eq:amplification} isolates the
same-day amplification branch. Equation~\eqref{eq:directed} corrects an early
mistaken duration interpretation, and Equation~\eqref{eq:pairsum} gives pair
balancing a precise invariant. These are model properties and implementation
contracts. No theorem yet bounds false-switch probability, expected detection
delay, or regret.

The exact loss-gap regularizer was not found in a limited primary-source scan,
but adjacent literature contains arbitrary transition matrices,
state-dependent persistence, continuous JMs, sparse feature selection, and
alternative selection criteria
\citep{aydinhan2024,nystrup2021,cortese2023}. The defensible phrase is
candidate mathematical extension, not a global priority claim.

\subsection{What has and has not been contributed economically}

The economic evidence is mixed. Germany is a large local development-sample
improvement: lagged Sharpe rises from 0.166440 to 0.337888, switches fall from
32 to 8, turnover falls from 0.993348 to 0.248337, and the strategy clears
both comparators. Pair balancing moves the US closer to HMM. Lagging improves
fixed JM Sharpe in Japan. But the final target is not met: the best observed US
JM still loses to HMM and the best observed Japanese JM still loses to
\BH{}.

This failure map is scientifically useful because it rejects a simple
universal story. A mechanism that suppresses candidate-path reversals can
still increase selected-strategy activity when monthly validation changes its
lambda schedule. The Japanese upper lambda is selected frequently, and the
path/choice interaction is large. Saying that the lagged path is simply better
is therefore too strong; lagging changes both path geometry and economic
selection, with heterogeneous net effects.

Even if a JM beat both benchmarks, it would demonstrate an economically useful
signal, not the existence of a unique true latent market state. The models are
unsupervised descriptions with no regime ground truth. Cash fraction also
matters: lower risk can arise mechanically from holding less equity. A valid
performance story requires improved return relative to risk, after equal
timing and cost, not simply more days in cash.

\subsection{Why the proxy differs from the parent paper}

The evidence supports underidentification rather than one confirmed bug:

\begin{enumerate}
\item Exact Bloomberg total-return indices and GFD cash series are unavailable.
The proxy test periods begin in 2007--2009 rather than 1990.
\item Japan uses a price index without dividends; German and Japanese cash
rates are different instruments from the parent Treasury bills.
\item The complete final JM and HMM candidate grids and several tie/calendar
conventions are undisclosed.
\item Grid choice is path-discrete and strongly binding. Restricting candidates
to Table-3-visible values, using a historical author grid, taking their union,
and adding one predefined endpoint all changed paths, but none restored the
three-market ordering.
\item The 5\% boundary diagnostic can label unresolved grid coverage, but it
cannot change any state or performance number.
\end{enumerate}

The endpoint audit added only JM \(362.0387\) and HMM smoothing \(1249\), both
derived before opening endpoint performance. The JM addition lowered US
Sharpe by 0.0739, raised German Sharpe by 0.0790 while worsening its MDD, and
left Japan unchanged. The HMM endpoint changed no primary metric. Grid
truncation is real, but this predefined test did not rescue the proxy.

\subsection{Limitations}

The current evidence has seven principal limits. First, proxy instruments and
test periods are not definition-equivalent to the parent study. Second, the
paper's complete grids remain unknown. Third, only three markets are used and
their crises are correlated. Fourth, adaptive centers were not jointly refit;
the experiments isolate decoding around frozen fixed-JM geometry. Fifth, beta
was scenario-fixed rather than estimated, and upper-lambda selection shows
that the finite adaptive optimum remains unidentified. Sixth, repeated
inspection through 2023 makes these markets development data. Seventh, no
confidence interval, multiple-testing correction, or untouched confirmation
sample supports a performance claim.

\section{Next mathematical hypothesis}

The current objective rewards feature compactness and temporal persistence.
It never requires a state to distinguish the future return earned by its
position. Market timing is therefore an indirect consequence. A simple next
principle is: a useful market state should be both statistically coherent and
economically different.

For training cutoff \(T\), a candidate predictive JM is
\begin{equation}
\min_{\Theta,\mu,S}
\sum_{u=0}^{T}\frac12\|x_u-\theta_{s_u}\|_2^2
+\gamma\sum_{u=0}^{T-2}
\rho\!\left(y_{u+2}-\mu_{s_u}\right)
+\lambda\sum_{u=1}^{T}\ind\{s_u\ne s_{u-1}\},
\label{eq:predictive}
\end{equation}
where \(y_{u+2}\) is the excess return matured two trading observations after
feature date \(u\), \(\mu_k\) is the return location associated with state
\(k\), and \(\rho\) is a robust prediction loss. The \(u+2\) label is legal
only if it is already observed by cutoff \(T\). At online date \(t\), inference
uses \(x_t\) and frozen training parameters; no future return enters the
decision.

Equation~\eqref{eq:predictive} changes the scientific direction from purely
unsupervised identification toward economically supervised state formation.
It has not been implemented or tested. Before an experiment, the relative
loss normalization, \(\gamma\), robust loss, alternating-fit algorithm,
state-label rule, nested selection, markets, and success criterion must be
frozen. Confirmation requires untouched markets or genuinely prospective
data. Further same-sample sweeps of robust-gap caps, two-day confirmation,
semi-Markov reversal surcharges, or feature metrics can generate hypotheses
but cannot manufacture confirmation from the inspected US/DE/JP sample.

\section{Conclusion}

Shu et al.\ begin with an investor, not an algorithm. Persistent unfavorable
markets create downside risk; a simple market/cash strategy makes the economic
value of a regime signal visible; HMM supplies a benchmark; and fixed JM asks
whether paying a constant cost for state changes produces a more useful,
tradable path. Their 1990--2023 evidence says yes.

We rebuilt that walk-forward question on public proxies and then asked whether every
switch should pay the same price. Arrival-day, lagged, and pair-balanced
transition penalties were formulated, implemented, and audited. The
mathematics is coherent: exact dynamic programming survives, beta zero is
fixed JM, same-day amplification has an explicit binary recursion, and pair
balancing has an exact identity.

The economic answer is not yet the desired one. The penalty formulas and
look-ahead-free accounting behaved as specified. Economic outcomes improved
clearly in one of three markets and were neutral or adverse in the other two.
Germany clears \BH{} and HMM; the US does not clear HMM; Japan
does not clear \BH{}. No all-market performance claim follows. The project has
advanced from an unexamined fixed barrier to a better-specified mathematical
family and an honest failure map. The next worthwhile test is not another
winner search on the same sample, but a frozen economically predictive state
objective evaluated on evidence that has not already shaped the model.

\appendix

\section{Exact candidate-grid inventory}

The following sets have distinct provenance and must not be conflated:

\begin{description}
\item[Paper-v3 Table-3 illustrations.]
JM \([0,5,15,35,70,150]\); HMM \([0,2,4,8,20]\).
They illustrate persistence and are not disclosed as final CV grids.

\item[Canonical proxy v7.]
JM \([0,5,15,35,70,150,300,600,1200]\);\\
HMM
\([0,2,4,6,8,10,20,40,80,160,320,640,1280,2560]\).

\item[Historical paper-v1 JM.]
\([10,22,50,100,220,500,1000]\), from an older and materially different
author design.

\item[Source-union audit.]
JM \([0,5,10,15,22,35,50,70,100,150,220,500,1000]\);\\
HMM
\([0,2,4,6,8,20]\).

\item[Performance-free behavior search.]
JM tested zero plus \(2^{j/2}\), \(j=-8,\ldots,22\); candidates through
\(362.0386719675\) were globally eligible and \(512+\) failed the frozen
occupancy/transition rule. HMM tested every integer \(0,\ldots,2560\);
\(0,\ldots,1249\) were globally eligible.

\item[Behavior-selected nine-candidate sets.]
\begingroup\small
JM \([0,0.353553,1,5.656854,16,32,64,181.019336,256]\);\\
HMM
\([0,3,9,32,54,114,166,402,1115]\).\\
The one-shot endpoint audit added only
JM \(362.038672\) and HMM \(1249\).
\endgroup
\end{description}

No tested grid recovers Shu et al.'s ordering across all three proxy markets.
The final paper grids remain unidentified. The adaptive experiments reuse the
canonical v7 lambda grid; they do not expand it.

\clearpage
\section{Chronology of mathematical and empirical development}

\begin{table}[!htbp]
\centering
\caption{Research sequence. Archived means the idea remains part of the
scientific record but is not current evidence.}
\label{tab:chronology}
\small
\begin{tabular}{p{0.13\textwidth}p{0.23\textwidth}p{0.53\textwidth}}
\toprule
Date/stage & Question or model & Outcome \\
\midrule
June 2026 & Minute prototype; hand-set additive/multiplicative costs and
\(\lambda=\log(d-1)\)
& Time-varying DP and synthetic paths worked, but real paths were 98--99.8\%
identical and transaction costs dominated. Loss/penalty scales were not
calibrated. Archived; no market evidence. \\
\addlinespace
8 July & Approximate daily P0 and exponential adaptive P1,
\(\lambda_t=\min\{\exp(b_0+b_1z_t),5000\}\)
& Superseded data, cadence, HMM, and accounting protocol. Historical Sharpe
deltas were negative in US/DE and slightly positive in JP. Archived null. \\
\addlinespace
8--10 July & Constant asymmetric exit/re-entry costs P2
& Binary identity~\eqref{eq:directed} disproved the intended independent
duration interpretation. Claim withdrawn; numbers describe a boundary bias. \\
\addlinespace
12 July & Causal fixed baseline v7
& Reproduced by independent replay. Fixed JM beat both benchmarks in 0/3
proxy markets. This is the accepted parent for adaptive work. \\
\addlinespace
14 July & 4,000-row JM fit window
& Upper-bound coverage failed in 8/9 market-delay cells, so performance stayed
sealed. Effect unknown. \\
\addlinespace
15--18 July & Table-3, historical, union, behavior-grid, and endpoint audits
& Grid choice materially changed paths, but no tested domain restored all
markets. The 5\% rule was confirmed descriptive only. \\
\addlinespace
17 July & Arrival evidence, adaptive-confidence-001
& Formula and causal tests passed; \(\beta=0\) exact. Trade-off result mixed,
but economic target passed 0/3. Same-day double-use motivated lagging. \\
\addlinespace
17 July & Training-prefix separation gate
& Leave-one-market-out diagnostic inconclusive and event ordering inconsistent.
Closed before P\&L. \\
\addlinespace
17--18 July & Lagged evidence mechanism and P\&L
& Candidate whipsaws \(17\to6\); selected Sharpe improved over fixed in all
markets, but only Germany beat both benchmarks. \\
\addlinespace
18 July & Path/choice attribution
& Large interaction; Shapley allocation assigned mean Sharpe gain mainly to
monthly choices, not uniformly better paths. Diagnostic, not causal inference. \\
\addlinespace
18--21 July & Pair-balanced mechanism and P\&L
& Pair identity exact and latency mostly retained, but whipsaws did not fall.
US improved; DE was nearly flat versus lagged; Japan worsened. Not supported. \\
\bottomrule
\end{tabular}
\end{table}

\clearpage
\section{Accepted experiment provenance}

The accepted evidence used in this manuscript is pinned by configs and
append-only registry entries:

\begin{itemize}
\item fixed baseline:
\path{fixed-baselines-8adb330565d6-3636939b525d-e9614112b234};
\item fixed assumption audit:
\path{fixed-baseline-assumption-audit-}\\
\path{79c94852c8fd-3636939b525d-4cc8cdbccd14};
\item arrival evidence:
\path{adaptive-confidence-1b0c327b2db4-3636939b525d-864d671cf973};
\item separation diagnostic:
\path{adaptive-separation-}\\
\path{813f66912526-26cbca8871be-fefc608b9081};
\item lagged mechanism:
\path{lagged-evidence-6f964f5724b2-26cbca8871be-d173ca32c86f};
\item lagged P\&L:
\path{lagged-pnl-bad599271e2d-643dd3e6d96f-be70588256b2};
\item lagged attribution:
\path{lagged-attribution-73a5995c487e-52854fc3c22a-197915169632};
\item endpoint audit:
\path{endpoint-grid-audit-05e9d08f619b-77b30ef98fa0-24ca06c297e8};
\item balanced mechanism:
\path{balanced-lagged-a7d9914ca1a8-643dd3e6d96f-17961bfd667f};
\item balanced P\&L:
\path{balanced-pnl-3ae665413a01-4e747110ba1c-eaae6444a9a5}.
\end{itemize}

Invalidated runs remain preserved for provenance but supply no accepted claim.
The final balanced numerical CSVs are byte-identical to two earlier runs
invalidated for provenance metadata only. The first understated the feature
columns physically loaded. The second labeled 43 explicitly locked files as
all files read even though inventory verification integrity-hashed 232 entries.
The accepted schema-2 lock separately records the 43 explicitly locked inputs,
all 232 integrity-hashed entries, all loaded feature columns, and the three
columns used for selection and accounting. Every scientific CSV, trade,
decision, smoke, and frozen lock is byte-identical across the corrected runs.

\section{Reproduction and availability}

The frozen sample cutoff is 31 December 2023. Repository source, lockfile,
configs, experiment registry, tests, and scientific ledger are versioned.
Raw provider downloads and generated run artifacts remain ignored. The core
commands are:

\begin{verbatim}
uv python install 3.12.3
uv sync --locked --extra data
.venv/bin/python -m pytest -q
.venv/bin/adaptive-jump fetch --config research.toml
.venv/bin/adaptive-jump run --study replication
  --config research.toml
.venv/bin/adaptive-jump verify
  --run artifacts/fixed-baselines/<run_id>
\end{verbatim}

The authors' \href{https://github.com/Yizhan-Oliver-Shu/jump-models}{public
jump-models repository} is a useful implementation reference, not the empirical
replication package for Shu et al.

\begin{thebibliography}{99}

\bibitem[Ang and Bekaert(2002)]{ang2002}
A. Ang and G. Bekaert.
\newblock International asset allocation with regime shifts.
\newblock \emph{Review of Financial Studies}, 15(4):1137--1187, 2002.

\bibitem[Ang and Bekaert(2004)]{ang2004}
A. Ang and G. Bekaert.
\newblock How regimes affect asset allocation.
\newblock \emph{Financial Analysts Journal}, 60(2):86--99, 2004.

\bibitem[Ang and Timmermann(2012)]{ang2012}
A. Ang and A. Timmermann.
\newblock Regime changes and financial markets.
\newblock \emph{Annual Review of Financial Economics}, 4(1):313--337, 2012.

\bibitem[Ayd{\i}nhan et al.(2024)]{aydinhan2024}
A. O. Ayd{\i}nhan, P. N. Kolm, J. M. Mulvey, and Y. Shu.
\newblock Identifying patterns in financial markets: Extending the statistical
jump model for regime identification.
\newblock \emph{Annals of Operations Research}, 2024.
\newblock doi:10.1007/s10479-024-06035-z.

\bibitem[Bemporad et al.(2018)]{bemporad2018}
A. Bemporad, V. Breschi, D. Piga, and S. P. Boyd.
\newblock Fitting jump models.
\newblock \emph{Automatica}, 96:11--21, 2018.

\bibitem[Bulla et al.(2011)]{bulla2011}
J. Bulla, S. Mergner, I. Bulla, A. Sesbo\"e, and C. Chesneau.
\newblock Markov-switching asset allocation: Do profitable strategies exist?
\newblock \emph{Journal of Asset Management}, 12:310--321, 2011.

\bibitem[Cortese et al.(2023)]{cortese2023}
F. P. Cortese, P. N. Kolm, and E. Lindstr\"om.
\newblock What drives cryptocurrency returns? A sparse statistical jump model
approach.
\newblock \emph{Digital Finance}, 5:483--518, 2023.

\bibitem[Hamilton(1989)]{hamilton1989}
J. D. Hamilton.
\newblock A new approach to the economic analysis of nonstationary time series
and the business cycle.
\newblock \emph{Econometrica}, 57(2):357--384, 1989.

\bibitem[Harding and Pagan(2003)]{harding2003}
D. Harding and A. Pagan.
\newblock A comparison of two business cycle dating methods.
\newblock \emph{Journal of Economic Dynamics and Control},
27(9):1681--1690, 2003.

\bibitem[Hardy(2001)]{hardy2001}
M. R. Hardy.
\newblock A regime-switching model of long-term stock returns.
\newblock \emph{North American Actuarial Journal}, 5(2):41--53, 2001.

\bibitem[Markowitz(1959)]{markowitz1959}
H. Markowitz.
\newblock \emph{Portfolio Selection}.
\newblock Yale University Press, New Haven, CT, 1959.

\bibitem[Nystrup et al.(2015)]{nystrup2015}
P. Nystrup, B. W. Hansen, H. Madsen, and E. Lindstr\"om.
\newblock Regime-based versus static asset allocation: Letting the data speak.
\newblock \emph{Journal of Portfolio Management}, 42(1):103--109, 2015.

\bibitem[Nystrup et al.(2018)]{nystrup2018}
P. Nystrup, B. W. Hansen, H. O. Larsen, H. Madsen, and E. Lindstr\"om.
\newblock Dynamic allocation or diversification: A regime-based approach to
multiple assets.
\newblock \emph{Journal of Portfolio Management}, 44(2):62--73, 2018.

\bibitem[Nystrup et al.(2020a)]{nystrup2020online}
P. Nystrup, P. N. Kolm, and E. Lindstr\"om.
\newblock Greedy online classification of persistent market states using
realized intraday volatility features.
\newblock \emph{Journal of Financial Data Science}, 2(3):25--39, 2020.

\bibitem[Nystrup et al.(2020b)]{nystrup2020persistent}
P. Nystrup, E. Lindstr\"om, and H. Madsen.
\newblock Learning hidden Markov models with persistent states by penalizing
jumps.
\newblock \emph{Expert Systems with Applications}, 150:113307, 2020.

\bibitem[Nystrup et al.(2021)]{nystrup2021}
P. Nystrup, P. N. Kolm, and E. Lindstr\"om.
\newblock Feature selection in jump models.
\newblock \emph{Expert Systems with Applications}, 184:115558, 2021.

\bibitem[Pagan and Sossounov(2003)]{pagan2003}
A. R. Pagan and K. A. Sossounov.
\newblock A simple framework for analysing bull and bear markets.
\newblock \emph{Journal of Applied Econometrics}, 18(1):23--46, 2003.

\bibitem[Roy(1952)]{roy1952}
A. D. Roy.
\newblock Safety first and the holding of assets.
\newblock \emph{Econometrica}, 20(3):431--449, 1952.

\bibitem[Ryd\'en et al.(1998)]{ryden1998}
T. Ryd\'en, T. Ter\"asvirta, and S. \AA sbrink.
\newblock Stylized facts of daily return series and the hidden Markov model.
\newblock \emph{Journal of Applied Econometrics}, 13(3):217--244, 1998.

\bibitem[Schwert(1989)]{schwert1989}
G. W. Schwert.
\newblock Why does stock market volatility change over time?
\newblock \emph{Journal of Finance}, 44(5):1115--1153, 1989.

\bibitem[Shu et al.(2024)]{shu2024}
Y. Shu, C. Yu, and J. M. Mulvey.
\newblock Downside risk reduction using regime-switching signals: A statistical
jump model approach.
\newblock \emph{Journal of Asset Management}, 25:493--507, 2024.
\newblock doi:10.1057/s41260-024-00376-x.

\bibitem[Smales(2016)]{smales2016}
L. Smales.
\newblock Risk-on/risk-off: Financial market response to investor fear.
\newblock \emph{Finance Research Letters}, 17:125--134, 2016.

\end{thebibliography}

\end{document}
